\section{Instrumenting it yourself}
\label{sec:ch03-instrumenting}
\index{diagnostics!execution listener|(}

Every measurement in this chapter came from three small files in Mosaic: a
listener, a per-request record for it to write into, and one options hook. They
are worth reading, partly because the timeline is genuinely useful to keep
switched on in development, and partly because getting them registered taught me
something that will bite again in Part~VI.

\subsection*{Scopes for phases, calls for facts}

\index{diagnostics!ExecutionDiagnosticEventListener}%
The extension point is \code{ExecutionDiagnosticEventListener}, a base class of
virtual no-ops that you subclass and override selectively. Its methods come in
two shapes, and the shape tells you what the event is for. A phase returns an
\code{IDisposable}: the engine opens the scope before the phase and disposes it
after, so timing a phase means timing the scope. A fact returns \code{void}.

\begin{minted}{csharp}
public override IDisposable ValidateDocument(RequestContext context)
    => Measure(context, static (timeline, elapsed) => timeline.Validate = elapsed);

public override void RetrievedDocumentFromCache(RequestContext context)
{
    if (Timeline(context) is { } timeline)
    {
        timeline.DocumentCacheHit = true;
    }
}
\end{minted}

\code{Measure} returns a small \code{IDisposable} holding a timestamp; there is
nothing to it beyond \code{Stopwatch.GetElapsedTime}. The two cache events are
facts rather than phases, which is how the timeline knows to print
\enquote{hit} or \enquote{miss}.

\index{diagnostics!EnableResolveFieldValue}%
Per-field events need one extra step. \code{ResolveFieldValue} fires once per
field rather than once per request, and the same switch gates \code{RunTask}, so
a listener has to ask for them:

\begin{minted}{csharp}
public override bool EnableResolveFieldValue => true;
\end{minted}

That property defaults to false, and the dispatcher reads it exactly once, in
its constructor, partitioning the listeners into those that want per-field
events and those that do not. So the cost of not opting in is zero rather than a
virtual call per field. In a service doing 146 resolver invocations per request
that is a defensible default, and it is worth knowing the switch exists before
you conclude your override is broken.

\subsection*{The registration that fails at startup}

\index{diagnostics!schema services}%
Registering the listener is one line, and next to it sits a line that looks like
boilerplate and is not:

\begin{minted}{csharp}
builder.AddGraphQL()
    .AddMosaic()
    .AddApplicationService<ILoggerFactory>()
    .AddDiagnosticEventListener<RequestTimelineListener>();
\end{minted}

\code{AddDiagnosticEventListener} registers the listener in the \emph{schema}
service provider, which is a separate container from the application's and does
not inherit its registrations. The listener wants an \code{ILoggerFactory}. The
application has one; the schema services do not, until you say so.

Delete that line and the service does not start:

\begin{minted}[fontsize=\footnotesize]{text}
fail: Microsoft.Extensions.Hosting.Internal.Host[11]
      Hosting failed to start
      System.InvalidOperationException: Unable to resolve service for type
      'Microsoft.Extensions.Logging.ILoggerFactory' while attempting to activate
      'Mosaic.Api.Infrastructure.Diagnostics.RequestTimelineListener'.
\end{minted}

I checked that by deleting it, because a chapter that tells you a line is
load-bearing should have watched the thing fall down. The message is a good one,
which is not something I get to say often in this book: it names both the
service it could not find and the type it was building, and it arrives at
startup rather than on the first query.

That last part is chapter~\ref{ch:hotchocolate-quickly}'s eager initialisation
paying out. Under the old lazy behaviour this would have surfaced when the first
request hit a listener that could not be constructed. Now the container never
becomes healthy and the deploy visibly does not happen. Error filters, HTTP
request interceptors and diagnostic listeners are all schema-level components
with this same requirement, and \code{AddApplicationService<T>} is the answer
for all of them.

\subsection*{Reading the pipeline back}

\index{request pipeline!reporting}%
The startup listing from section~\ref{sec:ch03-pipeline} comes from somewhere
else, because the pipeline is not something a diagnostic listener can see. It is
configured through named options, so it can be read back the same way:

\begin{minted}{csharp}
public static IServiceCollection AddPipelineReport(this IServiceCollection services)
    => services.AddTransient<IConfigureOptions<RequestExecutorSetup>>(
        serviceProvider => new ConfigureNamedOptions<RequestExecutorSetup>(
            ISchemaDefinition.DefaultName,
            setup => setup.PipelineModifiers.Add(
                pipeline => Report(serviceProvider, pipeline))));
\end{minted}

\code{Report} walks the list and logs each key. No reflection is involved: this
is the same mechanism \code{UseRequest} uses, reached through
\code{Microsoft.Extensions.Options} rather than through the internal helper
HotChocolate calls. The default schema is named \code{\_Default}, which is what
\code{ISchemaDefinition.DefaultName} holds.

The modifier only observes; it inserts nothing and removes nothing. Its position
in the registration order, as section~\ref{sec:ch03-pipeline} admitted, decides
whether it sees the finished list.

\subsection*{Verify it in Postman}

\index{Postman!pipeline requests}%
Two of the collection's requests exist to make the pipeline visible from the
client side, and both were added for this chapter.

The first sends the catalog-with-reviews query with \code{GraphQL-Cost:
validate}. It asserts that the response has no \code{data} property, and that
the field cost reported under \code{extensions.operationCost} is 30. That is a
client watching a request leave the pipeline at position seven.

The second posts \code{\{ products \{ title} with the brace left open and
asserts a 400 with \code{extensions.code} of \code{HC0011}. Run the collection
with the service's console visible and the point of
section~\ref{sec:ch03-parsing} lands better than any diagram: six of the seven
requests print a timeline line and that one does not.

\begin{minted}{text}
npx newman run postman/mosaic.postman_collection.json \
    -e postman/mosaic.local.postman_environment.json
\end{minted}

Seven requests, twenty-four assertions. \code{scripts/verify.ps1} runs the same
collection and also asserts the thirteen middleware in order and the 146
resolvers, so the numbers printed in this chapter are checked by the same script
that checks the numbers in the last one. If HotChocolate changes its default
pipeline, that script fails and this chapter is what needs rewriting.
\index{diagnostics!execution listener|)}
